\section{Introduction}

\subsection*{Abstract}

An introduction to partitions, with respect to discrete mathematics, as taught to me in my seminar sessions during my first year at Durham University. The following paper is as it was when it was originally submitted, with section one appended to the beginning. 

The paper sets out the basic terminology of partitions, before diving into some of the applications and specific combinatorics problems that can be calculated using them. 

\subsection*{Acknowledgements}

\thanks{In memory of the module professor, Professor Mikhail (Misha) Menshiokv, who passed away before the completion of this report.}

\section{Introduction to Partitions}

\subsection{Basic Definitions}

\definition{partition}{A \underline{partition} of \(n \in \N\) is a way of writing \(n\) as the sum of numbers. These will be denoted in this document as \((p_1 + p_2 + ... + p_i) \vdash n\).}

\separate{5}\definition{part}{A \underline{part} is one of the numbers that make up a partition. It has an associated \underline{part size}, which is the value of the number.}

\separate{5}\definition{partition function}{The \underline{partition function} of \(n\) is the number of possible partitions of \(n\). This will be denoted as \(p(n) = i\).}

\separate{5}\definition{set of partitions}{The \underline{set of partitions} is the set containing all possible partitions of \(n\). This will be denoted using non-standard notation as \(P(n) = \{p_1, p_2, ..., p_i\}\) where \(p_x\) is a partition of \(n\) and \(i = p(n)\).}

\separate{5}\definition{conditional partitions}{\underline{Conditional partitions} can also be created which limit which partitions are valid based on a given condition. This can be applied to both partition functions and sets of partitions as \(p(n|condition)\) and \(P(n|condition)\).}

\subsection{Bijective Proofs}

\definition{bijection}{A \underline{bijection} is a mathematical function that is a one-to-one and onto mapping \hyperref[cit:2]{(Merriam-Webster, 2025)}.} 

\section{Visualising Partitions}

\section{Generating Functions}

\section{Computing \textit{p(n)}}